% !TeX program = pdflatex
% Compilar: latexmk -pdf relatorio.tex   (ou pdflatex relatorio.tex duas vezes)
\documentclass[11pt]{article}
\usepackage[brazil]{babel}
\usepackage{relatorio}
\hypersetup{pdftitle={Pipeline experimental llm-ada-project}, pdfauthor={perrijuan}}

\title{Pipeline experimental llm-ada-project}
\author{Zero-shot, 3-shot, 5-shot e fine-tuning (LoRA) em quatro LLMs médicos}
\date{Documentação técnica do repositório -- outubro de 2026}
\paperurl{https://github.com/perrijuan/llm-ada-project}

\begin{document}
\thispagestyle{plain}
\maketitle

\begin{bloco}[Resumo]
O repositório compara quatro LLMs abertos (Gemma~4 E4B, HuatuoGPT-o1-7B, MedGemma~1.5~4B e Qwen3.5-4B) em quatro condições:
\textbf{zero-shot}, \textbf{3-shot}, \textbf{5-shot} (demonstrações no prompt, sem treino) e \textbf{fine-tuned} (SFT com adaptadores LoRA sobre o checkpoint em 4 bits).
A avaliação usa as mesmas 48 perguntas MedQuAD, 100 questões MedQA e 100 PubMedQA, além de um benchmark de velocidade (TTFT, TTLT, tokens/s).
\cod{rodar\_todos.py} executa um processo isolado por execução, \cod{registro.py} grava logs e tabelas, e \cod{consolidar.py} e \cod{resumir.py} produzem as tabelas finais.
Este documento descreve o código da versão atual (protocolo \cod{icl\_0\_3\_5\_ft5\_personas\_v4\_ft\_separado}).
\end{bloco}

\tableofcontents

% ------------------------------------------------------------------
\section{Visão geral do pipeline}

O experimento tem três camadas. O \textbf{executor} (\cod{rodar\_todos.py}) decide a ordem das execuções e abre um processo Python novo para cada uma.
O \textbf{script do modelo} (\cod{modelos/<modelo>.py}) prepara os dados, avalia as quatro condições e grava resultados.
As \textbf{ferramentas de pós-processamento} (\cod{consolidar.py} e \cod{resumir.py}) juntam as execuções concluídas e calculam os resumos estatísticos.

\begin{figure}[H]
\centering
\begin{tikzpicture}[node distance=6mm]
\node[caixa, text width=28mm] (exec) {\sbseries rodar\_todos.py\\[2pt]\footnotesize plano, trava,\\status, retomada};
\node[caixa, text width=28mm, right=8mm of exec] (mod) {\sbseries modelos/*.py\\[2pt]\footnotesize 1 processo\\por execução};
\node[caixa, text width=28mm, right=8mm of mod] (reg) {\sbseries registro.py\\[2pt]\footnotesize CSV, JSON, logs,\\figuras};
\node[caixa, text width=28mm, below=8mm of reg] (cons) {\sbseries consolidar.py\\[2pt]\footnotesize junta execuções\\com status ok};
\node[caixa, text width=28mm, left=8mm of cons] (res) {\sbseries resumir.py\\[2pt]\footnotesize média, IC 95\%,\\deltas pareados};
\node[caixa, text width=32mm, left=6mm of res] (pm) {\sbseries prompts\_medicos.py\\[2pt]\footnotesize persona, seeds\\finetuning\_dados.py};
\draw[seta] (exec) -- node[above, font=\scriptsize, text=DarkGray]{env} (mod);
\draw[seta] (mod) -- (reg);
\draw[seta] (reg) -- (cons);
\draw[seta] (cons) -- (res);
\draw[seta] (pm) -- (mod);
\draw[seta, dashed] (exec.south) |- ++(0,-4mm) -| node[pos=0.25, below, font=\scriptsize, text=DarkGray]{após cada processo} (cons.north west);
\end{tikzpicture}
\caption{Fluxo entre os arquivos principais. A consolidação é chamada pelo executor após cada processo e pode ser repetida manualmente.}
\end{figure}

Dentro de cada processo, o script do modelo segue sempre a mesma sequência:

\begin{figure}[H]
\centering
\begin{tikzpicture}
\node[caixa, text width=23mm] (a) {\sbseries Dados\\[2pt]\footnotesize filtro, dedup,\\split 192/48};
\node[caixa, text width=23mm, right=5mm of a] (b) {\sbseries Zero-shot\\[2pt]\footnotesize checkpoint\\original};
\node[caixa, text width=23mm, right=5mm of b] (c) {\sbseries 3-shot / 5-shot\\[2pt]\footnotesize recarrega,\\exemplos no prompt};
\node[caixa, text width=23mm, right=5mm of c] (d) {\sbseries Fine-tuned\\[2pt]\footnotesize recarrega, LoRA,\\5 épocas};
\node[caixa, text width=23mm, right=5mm of d] (e) {\sbseries Gráficos\\[2pt]\footnotesize e tabelas\\finais};
\draw[seta] (a) -- (b); \draw[seta] (b) -- (c); \draw[seta] (c) -- (d); \draw[seta] (d) -- (e);
\end{tikzpicture}
\caption{Sequência em um script de modelo. Em cada condição: MedQuAD (respostas, F1, ROUGE-L, NLL), velocidade, MedQA e PubMedQA.}
\end{figure}

A ordem das condições é apenas operacional. Cada condição recarrega o checkpoint original com \cod{build\_model()}, então nenhuma condição herda pesos ou estado da anterior.

% ------------------------------------------------------------------
\section{Modelos e condições experimentais}

\subsection{Os quatro modelos}

Os quatro scripts em \cod{modelos/} são praticamente idênticos: mudam o nome curto, o checkpoint, a pasta de saída e, no HuatuoGPT, o tratamento da resposta de raciocínio.

\begin{table}[H]
\centering\small
\footnotesize\begin{tabularx}{\textwidth}{@{}l l l X@{}}
\toprule
Chave & \cod{MODEL\_ID} & \cod{MODEL\_NAME} & Particularidade \\
\midrule
\cod{gemma4} & \cod{unsloth/gemma-4-E4B-it} & Gemma 4 E4B & Script de referência. \\
\cod{huatuogpt} & \cod{FreedomIntelligence/HuatuoGPT-o1-7B} & HuatuoGPT-o1-7B & Extrai o texto após \cod{\#\# Final Response}; previews com até 1.024 tokens. \\
\cod{medgemma} & \cod{unsloth/medgemma-1.5-4b-it} & MedGemma 1.5 4B & Igual ao Gemma, salvo o checkpoint. \\
\cod{qwen35} & \cod{unsloth/Qwen3.5-4B} & Qwen3.5-4B & Raciocínio desligado no chat template. \\
\bottomrule
\end{tabularx}
\caption{Modelos avaliados, na ordem padrão de execução.}
\end{table}

Todos carregam com \cod{FastModel.from\_pretrained} (Unsloth), em 4 bits (\cod{LOAD\_IN\_4BIT=True}) e com contexto de inferência \cod{MODEL\_CONTEXT\_LENGTH=16384}.
O chat template é sempre o do próprio tokenizer; \cod{enable\_thinking=False} só tem efeito nos templates que aceitam essa opção.

\subsection{As quatro condições}

As condições são definidas por uma única lista, lida como \textit{(nome, número de demonstrações, treina?)}:

\begin{lstlisting}[language=Python]
CONDITIONS = [('zero-shot', 0, False), ('3-shot', 3, False),
              ('5-shot', 5, False), ('fine-tuned', 0, True)]
\end{lstlisting}

\begin{table}[H]
\centering\small
\begin{tabularx}{\textwidth}{@{}l c c X@{}}
\toprule
Condição & Demonstrações & Treina pesos & O que muda \\
\midrule
zero-shot & 0 & não & Apenas \cod{system} (persona + instrução) e a pergunta. \\
3-shot & 3 & não & Três pares pergunta/resposta antes da pergunta de teste. \\
5-shot & 5 & não & Os mesmos três do 3-shot, mais dois, na mesma ordem. \\
fine-tuned & 0 & sim (LoRA) & SFT em 192 pares MedQuAD; avaliado sem demonstrações. \\
\bottomrule
\end{tabularx}
\caption{Condições experimentais. ``Shots'' contam demonstrações resolvidas, não personas.}
\end{table}

Few-shot (ICL) e fine-tuning são fluxos separados: as demonstrações nunca entram no \cod{SFTTrainer}, e o formatador de SFT (\cod{format\_training\_example}) não chama a seleção de demonstrações.

% ------------------------------------------------------------------
\section{Protocolo experimental}

\subsection{Dados e splits}

\begin{itemize}
\item \textbf{MedQuAD} (\cod{keivalya/MedQuad-MedicalQnADataset}): filtro \cod{keep\_example} (pergunta $\geq$ 12 caracteres, resposta entre 40 e 900 caracteres, sem prefixos genéricos como ``For more information\ldots''), deduplicação por pergunta normalizada (espaços e caixa), sorteio de 240 registros e split \textbf{192 treino / 48 teste} com \cod{split\_seed}.
\item \textbf{Banco de demonstrações MedQuAD}: \cod{indices\_demonstracoes} (em \cod{finetuning\_dados.py}) só aceita perguntas filtradas que não estão nem no treino nem no teste; exige pelo menos 5.
\item \textbf{MedQA} (\cod{GBaker/MedQA-USMLE-4-options}): 100 questões do split \cod{test}; demonstrações do split \cod{train}, sem perguntas do teste.
\item \textbf{PubMedQA} (\cod{qiaojin/PubMedQA}, \cod{pqa\_labeled}): 100 questões para teste; demonstrações dos registros restantes. Respostas \textit{yes/no/maybe} viram A/B/C.
\item \textbf{Velocidade}: as 20 primeiras perguntas do teste MedQuAD (\cod{SPEED\_N=20}), com exatamente 160 tokens gerados.
\item Os conjuntos usados ficam em \cod{registros/dados\_*.csv}, o que permite auditar a separação treino/teste/demonstrações.
\end{itemize}

\subsection{Seeds}

Há dois tipos de seed, separados de propósito:

\begin{table}[H]
\centering\small
\begin{tabularx}{\textwidth}{@{}l l X@{}}
\toprule
Seed & Opção & Controla \\
\midrule
\cod{split\_seed} & \cod{-{}-split-seed 42} & Subconjunto de 240, split 192/48, perguntas MedQA/PubMedQA. Fica fixa. \\
\cod{seed} & \cod{-{}-seeds 42 123 2026} & Seleção das demonstrações, \cod{set\_seed} antes de cada condição, \cod{random\_state} do LoRA e seed do \cod{SFTConfig}. \\
\bottomrule
\end{tabularx}
\caption{Seeds. Elas chegam ao script pelas variáveis \cod{LLM\_SPLIT\_SEED} e \cod{LLM\_SEED}.}
\end{table}

A seleção de demonstrações é determinística por pergunta: um gerador é criado a partir de \cod{sha256(seed|tarefa|pergunta)}, sorteia 5 exemplos do banco, e o 3-shot usa os 3 primeiros.
Assim, a mesma pergunta recebe os mesmos exemplos em todas as condições e modelos com a mesma seed; o 5-shot difere do 3-shot apenas pelos dois exemplos acrescentados.

\begin{lstlisting}[language=Python]
digest = hashlib.sha256(f'{SEED}|{task}|{question}'.encode('utf-8')).digest()
rng = random.Random(int.from_bytes(digest[:8], 'big'))
return rng.sample(demo_pools[task], 5)[:shots]
\end{lstlisting}

A variação entre seeds mistura ICL e aleatoriedade do treino, mas não mede variação do split. Repetições com a mesma seed (\cod{-{}-repeticoes}) medem repetibilidade, não amostras independentes, e não há garantia de determinismo numérico em todas as GPUs.

\subsection{Personas e estrutura das mensagens}

\cod{prompts\_medicos.py} define três personas (textos completos em \cod{prompts\_prontos.json}; fontes em \cod{PROMPTS\_E\_PROTOCOLO.md}):

\begin{itemize}
\item \cod{clinico} (padrão): médico generalista, adaptado do MedAgents.
\item \cod{biomedico}: pesquisador biomédico, adaptado da model card OpenBioLLM.
\item \cod{original}: instruções do pacote recebido, sem persona.
\end{itemize}

O prompt de sistema é \textit{prefixo da persona + instrução da tarefa}: ``Answer the medical question directly in 2-4 factual sentences.'' no MedQuAD e ``Return only the option letter\ldots'' em MedQA/PubMedQA.
A persona é a mesma nas quatro condições de uma campanha; trocar a persona exige outra pasta de saída.

\begin{lstlisting}
system:    persona + instrucao da tarefa
user:      pergunta de demonstracao 1
assistant: resposta de referencia 1
...        (3 ou 5 pares; nenhum no zero-shot e no fine-tuned)
user:      pergunta de teste (sem resposta)
\end{lstlisting}

\subsection{Fine-tuning (SFT com LoRA)}

O braço \textit{fine-tuned} recarrega o checkpoint base, aplica \cod{FastModel.get\_peft\_model} e treina com \cod{SFTTrainer} (TRL). Cada exemplo de treino é uma conversa \textit{system + pergunta + resposta de referência}.

\begin{table}[H]
\centering\small
\begin{tabular}{@{}l l@{\hspace{6mm}}l l@{}}
\toprule
Parâmetro & Valor & Parâmetro & Valor \\
\midrule
Exemplos de treino & 192 & Épocas & 5 (\cod{max\_steps=-1}) \\
Quantização base & 4 bits (QLoRA) & Learning rate & $2\times10^{-5}$, linear \\
LoRA $r$ / $\alpha$ / dropout & 8 / 16 / 0 & Warm-up & 5 passos \\
Módulos & atenção + MLP (linguagem) & Otimizador & \cod{adamw\_torch}, wd 0{,}01 \\
Microbatch $\times$ acumulação & $2\times16=32$ & Comprimento SFT & 768 tokens \\
Precisão & bf16 se suportado, senão fp16 & Checkpointing & \cod{'unsloth'} \\
\bottomrule
\end{tabular}
\caption{Hiperparâmetros de \cod{train\_regime()}. O adapter e o tokenizer são salvos em \cod{<OUTPUT\_DIR>/adapter-fine-tuned}.}
\end{table}

\subsection{Guardrails}

O código prefere falhar com mensagem clara a mudar o protocolo em silêncio. As principais verificações:

\begin{table}[H]
\centering\small
\begin{tabularx}{\textwidth}{@{}l X@{}}
\toprule
Onde & Verificação \\
\midrule
script do modelo & Exige CUDA e \cod{WORLD\_SIZE=1} (uma GPU, para manter o batch efetivo 32). \\
\cod{select\_demonstrations} & \cod{shots} precisa ser 0, 3 ou 5. \\
\cod{make\_messages} & Fine-tuned nunca recebe demonstrações. \\
\cod{encode\_prompt} & Prompt + tokens reservados para a resposta não podem passar de 16.384; não corta exemplos. \\
\cod{make\_demo\_pool} / \cod{indices\_demonstracoes} & Pelo menos 5 demonstrações sem sobreposição com teste (e, no MedQuAD, com treino). \\
\cod{textos\_para\_tokenizacao} & Textos de SFT precisam ser \cod{str} não vazias; indica o índice inválido. \\
\cod{answer\_nll} & O prefixo tokenizado da conversa completa deve ser igual ao prompt; resposta precisa ter tokens. \\
\cod{TokenTimer} / \cod{timed\_generate} & Batch 1, um token por callback, tokens do streamer iguais aos do \cod{generate}, 160 tokens exatos no benchmark. \\
\cod{rodar\_todos.py} & Seeds no intervalo $[0, 2^{32})$, sem modelos ou seeds repetidos, \cod{-{}-repeticoes} $\geq 1$. \\
\bottomrule
\end{tabularx}
\caption{Verificações que interrompem a execução com \cod{ValueError} ou \cod{RuntimeError}.}
\end{table}

% ------------------------------------------------------------------
\section{Orquestração com rodar\_todos.py}

\subsection{Opções}

\begin{table}[H]
\centering\small
\begin{tabularx}{\textwidth}{@{}l l X@{}}
\toprule
Opção & Padrão & Efeito \\
\midrule
\cod{-{}-modelos} & os 4 & Quais modelos e em que ordem (\cod{gemma4 huatuogpt medgemma qwen35}). \\
\cod{-{}-seeds} & \cod{42} & Seeds de execução. \\
\cod{-{}-repeticoes} & \cod{10} & Repetições por seed (use \cod{1} para estudar seeds). \\
\cod{-{}-split-seed} & \cod{42} & Seed fixa dos dados. \\
\cod{-{}-persona} & \cod{clinico} & \cod{clinico}, \cod{biomedico} ou \cod{original}. \\
\cod{-{}-saida} & \cod{resultados\_icl\_personas} & Pasta da campanha. \\
\cod{-{}-planejar} & -- & Só imprime o plano; não carrega modelos, dados nem GPU. \\
\cod{-{}-retomar} & -- & Pula execuções completas com mesmo código e configuração. \\
\bottomrule
\end{tabularx}
\caption{Argumentos de linha de comando.}
\end{table}

\subsection{Ordem e isolamento}

O executor monta a lista de trabalhos com três laços: \textbf{modelo} $\rightarrow$ \textbf{seed} $\rightarrow$ \textbf{repetição}. Assim, um modelo termina todas as suas execuções antes de o próximo começar.
Com \cod{-{}-seeds 42 123 2026 -{}-repeticoes 1} são 12 processos e 12 fine-tunings.

Para cada trabalho, o executor cria \cod{<saida>/<modelo>/run\_NN/}, grava \cod{status.json} com \cod{status=\textquotedbl{}rodando\textquotedbl{}} e chama o script com \cod{subprocess.Popen}, usando a pasta como diretório de trabalho. A configuração vai por variáveis de ambiente:

\begin{lstlisting}
LLM_RUN_DIR, LLM_SEED, LLM_SPLIT_SEED, LLM_PERSONA,
PYTHONUNBUFFERED=1, PYTHONIOENCODING=utf-8, MPLBACKEND=Agg
\end{lstlisting}

A saída do processo é mostrada no terminal e copiada linha a linha para \cod{execucao.log}. O executor espera o processo terminar antes de abrir o próximo.

\subsection{Estados, falhas e retomada}

\begin{figure}[H]
\centering
\begin{tikzpicture}
\node[caixa, text width=22mm] (r) {\sbseries rodando};
\node[caixa, text width=30mm, right=14mm of r, yshift=14mm] (ok) {\sbseries ok\\\footnotesize código 0};
\node[caixa, text width=30mm, right=14mm of r] (er) {\sbseries erro\\\footnotesize código $\neq 0$ ou exceção};
\node[caixa, text width=30mm, right=14mm of r, yshift=-14mm] (in) {\sbseries interrompido\\\footnotesize Ctrl+C};
\draw[seta] (r.east) -- (ok.west); \draw[seta] (r.east) -- (er.west); \draw[seta] (r.east) -- (in.west);
\node[right=6mm of ok, text width=60mm, font=\footnotesize, align=left] {entra na consolidação};
\node[right=6mm of er, text width=60mm, font=\footnotesize, align=left] {pula as demais execuções do modelo e segue para o próximo; código final 1};
\node[right=6mm of in, text width=60mm, font=\footnotesize, align=left] {encerra o processo filho e a campanha};
\end{tikzpicture}
\caption{Estados gravados em \cod{status.json} (com \cod{duracao\_s} e \cod{exit\_code}).}
\end{figure}

\begin{itemize}
\item \textbf{Falha}: o traceback vai para \cod{execucao.log}, o processo filho é encerrado (\cod{terminate}, depois \cod{kill} após 20\,s), resultados parciais ficam na pasta e são excluídos da consolidação.
\item \textbf{Trava}: o arquivo \cod{.em\_execucao} impede dois executores na mesma campanha.
\item \textbf{Manifesto}: \cod{campanha.json} fixa protocolo, seeds, split, persona e repetições; se a configuração mudar, o executor pede outra \cod{-{}-saida}. Outro protocolo na mesma pasta também é recusado.
\item \textbf{Retomada}: com \cod{-{}-retomar}, uma pasta existente só é pulada se \cod{status=\textquotedbl{}ok\textquotedbl{}} e todos os campos esperados batem, inclusive \cod{codigo\_sha256} (hash do script do modelo e de todos os \cod{*.py} da raiz). Pastas com erro não são sobrescritas; não se retoma uma época no meio.
\end{itemize}

% ------------------------------------------------------------------
\section{Artefatos e logs}

\subsection{Estrutura de uma campanha}

\begin{lstlisting}
campanha_clinico/
  campanha.json                 configuracao fixa da campanha
  consolidado/                  saida de consolidar.py e resumir.py
  gemma4/run_01/
    status.json  execucao.log  ambiente.json  nvidia-smi.txt
    registros/                  prompts, dados, treino, tabelas
    figuras/                    PNGs dos graficos
    gemma4-e4b-medquad-unsloth/adapter-fine-tuned/
    gemma4-e4b-medquad-unsloth-extended/   MedQuAD por pergunta e resumo
    gemma4-e4b-medquad-unsloth-benchmark/  velocidade, MedQA, PubMedQA
  huatuogpt/run_01/ ...
\end{lstlisting}

\subsection{registro.py}

\cod{registro.py} só registra; não altera pesos, loss ou hiperparâmetros.

\begin{table}[H]
\centering\small
\begin{tabularx}{\textwidth}{@{}l X@{}}
\toprule
Função & O que faz \\
\midrule
\cod{escrever\_csv} & Escreve em \cod{.tmp} e renomeia (escrita atômica), UTF-8 com BOM, união das colunas de todas as linhas; listas e dicionários viram JSON. \\
\cod{iniciar} & Cria pasta própria se o script rodar sozinho; grava \cod{ambiente.json} (Python, plataforma, versões de torch/unsloth/transformers/trl/peft\ldots, persona, prompts, seeds, hiperparâmetros) e \cod{nvidia-smi.txt}. \\
\cod{etapa} & Imprime \cod{[ETAPA n] descrição} no log. \\
\cod{registrar\_trainer} & Callback que grava \cod{treino\_<tag>\_steps.csv} (log\_history) e \cod{config\_<tag>.json} (argumentos resolvidos do Trainer). \\
\cod{salvar\_resultado} & \cod{treino\_<tag>\_metricas.csv} com \cod{training\_loss}, \cod{global\_step} e métricas finais. \\
\cod{salvar\_tabelas} & Linhas e DataFrames do script, conjuntos \cod{dados\_*.csv}, \cod{datasets.csv} (tamanho e fingerprint) e \cod{model\_config.json}. \\
\cod{salvar\_figuras} & Salva todas as figuras abertas em \cod{figuras/} (130 dpi). \\
\bottomrule
\end{tabularx}
\caption{Funções de \cod{registro.py}.}
\end{table}

\begin{table}[H]
\centering\small
\begin{tabularx}{\textwidth}{@{}l X@{}}
\toprule
Arquivo & Conteúdo \\
\midrule
\cod{benchmark\_summary.csv} & Velocidade (média, mediana, P99/P1), acurácia com IC e métricas macro, por condição. \\
\cod{speed\_per\_question.csv} & TTFT, TTLT, tempo total, tokens/s, \cod{prompt\_tokens}, \cod{token\_times}. \\
\cod{precision\_per\_question.csv} & Letra prevista, resposta e acerto em MedQA/PubMedQA. \\
\cod{answers\_per\_question.csv} & Resposta MedQuAD, referência, tempos, F1, ROUGE-L, NLL. \\
\cod{metrics\_summary.csv} & Agregado MedQuAD e métricas de treino por condição. \\
\cod{registros/prompts.csv} & Mensagens completas, IDs das demonstrações, seed, persona, condição. \\
\cod{registros/previews.csv} & Três respostas completas por condição. \\
\cod{registros/treino\_fine-tuned\_*} & Loss, learning rate e métricas finais do SFT. \\
\bottomrule
\end{tabularx}
\caption{Principais arquivos de resultado por execução.}
\end{table}

\subsection{Pasta auditoria/}

\cod{auditoria/} guarda o histórico das revisões do pacote: \cod{manifesto.json} e \cod{manifesto\_personas\_v3.json} (SHA-256 dos notebooks, scripts e documentos entregues), os diffs de cada script em relação ao original (\cod{gemma4.diff}, \cod{personas\_v3.diff}, \cod{correcao\_finetuning.diff}\ldots), \cod{PLANO\_PERSONAS\_V3.txt} (saída de \cod{-{}-planejar}) e as notas \cod{VALIDACAO.md}, \cod{REVISAO\_INICIAL\_HISTORICA.md} e \cod{REVISAO\_PERSONAS\_V3.md}.
Esses arquivos descrevem versões anteriores; o código atual e o \cod{README.md} prevalecem.

% ------------------------------------------------------------------
\section{Consolidação e métricas}

\subsection{Métricas por execução}

\textbf{MedQuAD (resposta aberta, até 160 tokens, greedy).}
Token F1 compara palavras normalizadas com multiplicidade; ROUGE-L usa a maior subsequência comum de palavras. Ambos medem semelhança lexical, não correção médica.
A NLL soma a perda dos tokens da resposta de referência (dado o prompt da condição) e divide pelo total de tokens:
\[
\mathrm{NLL} = \frac{\sum_i \mathrm{nll}_i}{\sum_i n_i}, \qquad \mathrm{PPL} = e^{\mathrm{NLL}} .
\]
Também são registrados \cod{finished\_rate} (resposta terminou com EOS), número de palavras e pico de VRAM alocada pelo PyTorch.

\textbf{Velocidade.} O \cod{TokenTimer} (um \cod{BaseStreamer}) registra o instante de cada token gerado, com sincronização CUDA antes e depois do \cod{generate}. Há um warm-up de 8 tokens fora da medição.
\[
\mathrm{TTFT} = t_1 - t_0, \qquad \mathrm{TTLT} = t_N - t_0, \qquad \text{tokens/s} = \frac{N-1}{\mathrm{TTLT}-\mathrm{TTFT}} .
\]
\cod{benchmark\_summary.csv} traz média, mediana e a cauda do pior caso: P99 para tempos (menor é melhor) e P1 para tokens/s (maior é melhor). Com $n=20$ o P99 é quase o máximo, e o código emite um aviso (use $n\geq100$).

\textbf{MedQA e PubMedQA.} \cod{pick\_option} acrescenta \cod{\textquotedbl{}Answer:\textquotedbl{}} ao prompt e escolhe a letra com o maior logit do próximo token (considerando ``A'' e `` A''). A acurácia tem IC aproximado $\bar{x} \pm 1{,}96\, s/\sqrt{n}$, limitado a $[0,1]$.
Precision, recall e F1 macro são calculados por alternativa (A--D ou A--C) e depois tirada a média simples; não são F1 por especialidade.

\subsection{consolidar.py}

\begin{itemize}
\item Lê todos os \cod{*/run\_*/status.json}; se houver mais de um protocolo na pasta, para com erro.
\item Grava \cod{consolidado/execucoes.csv} com todas as execuções, inclusive as que falharam.
\item Para execuções com \cod{status=\textquotedbl{}ok\textquotedbl{}}, junta cada CSV pelo nome do arquivo, ignora \cod{dados\_*} e acrescenta \cod{exec\_modelo}, \cod{exec\_repeticao}, \cod{exec\_seed}, \cod{exec\_split\_seed}, \cod{exec\_persona} e \cod{exec\_origem}.
\item CSVs da pasta \cod{registros/} recebem o prefixo \cod{registro\_\_}; são cópias e não devem ser somadas aos CSVs originais.
\end{itemize}

\subsection{resumir.py}

Agrega entre execuções. Para cada grupo calcula $n$, média, desvio-padrão, mínimo, máximo e IC 95\% da média com $t$ de Student:
\[
\bar{x} \pm t_{0{,}975,\,n-1}\,\frac{s}{\sqrt{n}} .
\]

\begin{table}[H]
\centering\small
\begin{tabularx}{\textwidth}{@{}l X@{}}
\toprule
Saída em \cod{consolidado/} & Conteúdo \\
\midrule
\cod{resumo\_entre\_execucoes.csv} & Resumo de \cod{benchmark\_summary} por modelo, categoria, métrica e condição. \\
\cod{deltas\_pareados.csv} & Diferença de cada condição contra o zero-shot da mesma execução. \\
\cod{resumo\_deltas.csv} & Resumo desses deltas (3-shot, 5-shot e fine-tuned $-$ zero-shot). \\
\cod{resumo\_medquad\_condicoes.csv} & Resumo de \cod{metrics\_summary} por modelo, regime e métrica. \\
\bottomrule
\end{tabularx}
\caption{Arquivos gerados por \cod{resumir.py}.}
\end{table}

O IC entre execuções descreve a variação observada entre execuções; não substitui a incerteza entre perguntas.

% ------------------------------------------------------------------
\section{Como reproduzir}

\textbf{Ambiente.} É preciso uma GPU NVIDIA com CUDA e Python com PyTorch, Unsloth, Transformers e TRL compatíveis com os checkpoints. \cod{requirements.txt} lista só as dependências auxiliares (\cod{tqdm}, \cod{pandas}, \cod{numpy}, \cod{matplotlib}, \cod{ipython}, \cod{datasets}, \cod{sentencepiece}, \cod{huggingface\_hub}, \cod{scipy}); a pilha de treino segue a instalação oficial do Unsloth. Checkpoints restritos exigem login no Hugging Face.

\begin{lstlisting}[language=bash]
# 1. Dependencias auxiliares (depois de PyTorch CUDA + Unsloth)
python -m pip install -r requirements.txt

# 2. Conferir o plano (sem GPU, sem downloads)
python rodar_todos.py --repeticoes 1 --seeds 42 123 2026 --persona clinico --planejar

# 3. Piloto com um modelo
python rodar_todos.py --modelos medgemma --repeticoes 1 --seeds 42 --saida piloto
python -m pip freeze > ambiente_validado.txt

# 4. Campanha completa: 4 modelos x 3 seeds = 12 processos
python rodar_todos.py --repeticoes 1 --seeds 42 123 2026 --persona clinico --saida campanha_clinico

# 5. Retomar apos interrupcao (pula so execucoes completas)
python rodar_todos.py --repeticoes 1 --seeds 42 123 2026 --persona clinico --saida campanha_clinico --retomar

# 6. Consolidar e resumir (a consolidacao tambem roda apos cada processo)
python consolidar.py campanha_clinico
python resumir.py campanha_clinico
\end{lstlisting}

\begin{bloco}[Boas práticas]
\begin{itemize}
\item Não atualize bibliotecas durante uma campanha; use a mesma GPU para todos os modelos.
\item Mudou código, persona, seeds ou repetições: use uma nova \cod{-{}-saida}.
\item Para estudar seeds, prefira \cod{-{}-repeticoes 1}; repetições da mesma seed não são amostras independentes.
\item Outra persona é outra campanha: \cod{-{}-persona biomedico -{}-saida campanha\_biomedico}.
\end{itemize}
\end{bloco}

% ------------------------------------------------------------------
\section{Limitações e pontos de atenção}

\begin{itemize}
\item \textbf{Sem execução real nesta documentação}: o texto foi escrito a partir do código; não houve download de modelos, inferência ou treino em GPU.
\item \textbf{Testes ausentes}: o \cod{README.md} cita \cod{python -m unittest discover -s testes -v}, mas a pasta \cod{testes/} foi removida do repositório; o comando não funciona na versão atual.
\item \textbf{Nome do protocolo}: \cod{status.json} e \cod{campanha.json} usam \cod{icl\_0\_3\_5\_ft5\_personas\_v4\_ft\_separado} (\cod{rodar\_todos.py}), mas \cod{ambiente.json} grava \cod{icl\_0\_3\_5\_ft5\_personas\_v3} (\cod{prompts\_medicos.py}). Não afeta a execução; vale alinhar para evitar confusão na auditoria.
\item \textbf{Retomada}: \cod{codigo\_sha256} cobre os \cod{*.py}, mas não \cod{prompts\_prontos.json} (que é só referência; os prompts usados vêm de \cod{prompts\_medicos.py}).
\item \textbf{HuatuoGPT}: avaliação principal e velocidade usam 160 tokens, o que pode cortar o raciocínio antes de \cod{\#\# Final Response}.
\item \textbf{Métricas}: F1 e ROUGE-L são lexicais; perplexidade depende do tokenizador e não serve como ranking direto entre famílias; a escolha de letra pelo próximo token supõe letras de um único token.
\item \textbf{Comparação ICL $\times$ FT}: em MedQA e PubMedQA o few-shot usa exemplos do próprio benchmark, enquanto o fine-tuning usa apenas MedQuAD.
\end{itemize}

\end{document}
